% META: {"id": "1NSI-AGT-CO-001", "chapitre": "1NSI-ALGO-PARCOURS-TRIS", "type_objet": "corrige", "exercice_ref": "1NSI-AGT-EX-001", "capacites": ["P-ALGO-01B"], "sources_inspiration": [], "status": "verified"}
% BEGIN-VERIFY
% def minimum(tableau):
%     mini = tableau[0]
%     for x in tableau[1:]:
%         if x < mini:
%             mini = x
%     return mini
% def moyenne(tableau):
%     return sum(tableau) / len(tableau)
% notes = [12, 4, 9, 1, 7]
% assert minimum(notes) == 1
% assert moyenne(notes) == 6.6
% END-VERIFY
\begin{corrige}{1NSI-AGT-EX-001}
\textbf{1.}
\begin{python}
def minimum(tableau):
    mini = tableau[0]
    for x in tableau[1:]:
        if x < mini:
            mini = x
    return mini
\end{python}

\textbf{2.} Le minimum est $1$. La moyenne est
$\dfrac{12+4+9+1+7}{5}=\dfrac{33}{5}=6{,}6$.

\textbf{3.} Ce choix fonctionne seulement parce que l'élève connaît une \textbf{borne
maximale a priori} (ici, $20$, la note maximale). Cette astuce n'est pas généralisable :
pour un tableau de données quelconques (températures, scores non bornés...), il n'existe
pas toujours de constante connue à l'avance garantie supérieure à tous les éléments.
Initialiser avec le \textbf{premier élément} du tableau fonctionne, elle, dans
\textbf{tous} les cas, sans connaissance préalable des données.
\end{corrige}
